\par\Needspace{6\baselineskip}\medskip
\hypertarget{native-159-article114}{}\pdfbookmark[2]{114.}{article-114}
\begin{center}114.\end{center}
Multiplicando aequationem 1 per $\sin\alpha''\operatorname{tang}B''-\sin L''\operatorname{tang}\beta''$, aequationem 2 per $\cos L''\operatorname{tang}\beta''-\cos\alpha''\operatorname{tang}B''$, aequationem 3 per $\sin(L''-\alpha'')$, addendoque producta, prodit
\[
\begin{aligned}
0={}&n\{(0.2.\mathrm{II})\delta+(\mathrm{O}.2.\mathrm{II})D\}\\
&-n'\{(1.2.\mathrm{II})\delta'+(\mathrm{I}.2.\mathrm{II})D'\}
\end{aligned}\tag*{4)}
\]
similique modo, vel commodius per solam locorum inter se permutationem
\[
\begin{aligned}
0={}&n\{(0.1.\mathrm{I})\delta+(\mathrm{O}.1.\mathrm{I})D\}\\
&+n''\{(2.1.\mathrm{I})\delta''+(\mathrm{II}.1.\mathrm{I})D''\}
\end{aligned}\tag*{5)}
\]
\[
\begin{aligned}
0={}&n'\{(1.0.\mathrm{O})\delta'+(\mathrm{I}.0.\mathrm{O})D'\}\\
&-n''\{(2.0.\mathrm{O})\delta''+(\mathrm{II}.0.\mathrm{O})D''\}.
\end{aligned}\tag*{6)}
\]
Quodsi itaque ratio quantitatum $n$, $n'$ data est, adiumento aequationis 4 ex $\delta$ determinare licebit $\delta'$, vel $\delta$ ex $\delta'$; similiterque de aequationibus 5, 6. E combinatione aequationum 4, 5, 6 oritur haec
\[
\begin{aligned}
&\frac{(0.2.\mathrm{II})\delta+(\mathrm{O}.2.\mathrm{II})D}{(0.1.\mathrm{I})\delta+(\mathrm{O}.1.\mathrm{I})D}\\[.6em]
&\qquad\times\frac{(1.0.\mathrm{O})\delta'+(\mathrm{I}.0.\mathrm{O})D'}{(1.2.\mathrm{II})\delta'+(\mathrm{I}.2.\mathrm{II})D'}\\[.6em]
&\qquad\times\frac{(2.1.\mathrm{I})\delta''+(\mathrm{II}.1.\mathrm{I})D''}{(2.0.\mathrm{O})\delta''+(\mathrm{II}.0.\mathrm{O})D''}=-1,
\end{aligned}\tag*{7)}
\]
per quam e duabus distantiis corporis coelestis a terra determinare licet tertiam. Ostendi potest autem, hanc aequationem 7 fieri identicam, adeoque ad determinationem unius distantiae e duabus reliquis ineptam, quoties fuerit $B=B'=B''=0$ atque
\[
\left.\begin{aligned}
&\operatorname{tang}\beta'\operatorname{tang}\beta''\sin(L-\alpha)\sin(L''-L')\\
&+\operatorname{tang}\beta''\operatorname{tang}\beta\sin(L'-\alpha')\sin(L-L'')\\
&+\operatorname{tang}\beta\operatorname{tang}\beta'\sin(L''-\alpha'')\sin(L'-L)
\end{aligned}\right\}=0.
\]
\par
Ab hoc incommodo libera est formula sequens, ex aequationibus 1, 2, 3 \SourceLeaf{160}facile demanans:
\[
\begin{aligned}
&(0.1.2)\delta\delta'\delta''+(\mathrm{O}.1.2)D\delta'\delta''\\
&+(0.\mathrm{I}.2)D'\delta\delta''+(0.1.\mathrm{II})D''\delta\delta'\\
&+(0.\mathrm{I}.\mathrm{II})D'D''\delta+(\mathrm{O}.1.\mathrm{II})DD''\delta'\\
&+(\mathrm{O}.\mathrm{I}.2)DD'\delta''+(\mathrm{O}.\mathrm{I}.\mathrm{II})DD'D''=0.
\end{aligned}\tag*{8)}
\]
\par
Multiplicando aequationem 1 per $\sin\alpha'\operatorname{tang}\beta''-\sin\alpha''\operatorname{tang}\beta'$, aequationem 2 per $\cos\alpha''\operatorname{tang}\beta'-\cos\alpha'\operatorname{tang}\beta''$, aequationem 3 per $\sin(\alpha''-\alpha')$, addendoque producta, prodit
\[
\begin{aligned}
0={}&n\{(0.1.2)\delta+(\mathrm{O}.1.2)D\}\\
&-n'(\mathrm{I}.1.2)D'+n''(\mathrm{II}.1.2)D''
\end{aligned}\tag*{9)}
\]
et perinde
\[
\begin{aligned}
0={}&n(0.\mathrm{O}.2)D-n'\{(0.1.2)\delta'+(0.\mathrm{I}.2)D'\}\\
&+n''(0.\mathrm{II}.2)D''
\end{aligned}\tag*{10)}
\]
\[
\begin{aligned}
0={}&n(0.1.\mathrm{O})D-n'(0.1.\mathrm{I})D'\\
&+n''\{(0.1.2)\delta''+(0.1.\mathrm{II})D''\}.
\end{aligned}\tag*{11)}
\]
Adiumento harum aequationum e ratione inter quantitates $n$, $n'$, $n''$ cognita eruere licebit distantias $\delta$, $\delta'$, $\delta''$. Sed haecce conclusio generaliter tantum loquendo valet, exceptionemque patitur, quoties fit $(0.1.2)=0$. Ostendi enim potest, in hocce casu ex aequationibus 9, 10, 11 nihil aliud sequi, nisi relationem necessariam inter quantitates $n$, $n'$, $n''$, et quidem e singulis tribus eandem. Restrictiones analogae circa aequationes 4, 5, 6 lectori perito sponte se offerent.
\par
Ceterum omnes conclusiones hic evolutae nullius sunt usus, quoties planum orbitae cum ecliptica coincidit. Si enim $\beta$, $\beta'$, $\beta''$, $B$, $B'$, $B''$ omnes sunt $=0$, aequatio 3 \emph{identica} est, ac proin omnes quoque sequentes.
\par\bigskip\begin{center}\rule{.23\linewidth}{.35pt}\end{center}
